


The choice of neural network architecture and training hyperparameters plays a crutial role in the performance of physics-informed neural networks.
loss landscape

The original PINN framework was demonstrated mainly on relatively low-dimensional problems, including examples such as the Burgers equation, the Schrödinger equation, and the Allen--Cahn equation \cite{pinns_original}. These early implementations typically used a fixed set of collocation points sampled before training, often by Latin hypercube sampling. The number of residual points was commonly on the order of hundreds to thousands, and the networks were comparatively small. For example, architectures such as nine hidden layers with twenty neurons per layer, five hidden layers with one hundred neurons per layer, or four hidden layers with two hundred neurons per layer were sufficient for many one- and two-dimensional benchmark problems. Optimization was frequently performed using the L-BFGS algorithm, either alone or after an initial stage of first-order optimization.

When moving toward higher-dimensional diffusion problems, the requirements on the network and the training procedure become more demanding. The solution must be represented over a larger input space, and the residual loss has to capture the behaviour of the differential operator in many directions. Increasing the network width is therefore a natural way to improve representational capacity. Recent studies on the optimization of PINNs have investigated widths such as 50, 100, 200, and 400 neurons per layer, using three hidden layers and approximately $10^4$ residual points \cite{loss_landscape}. Such experiments often rely on Glorot initialization and compare different training lengths, for example $10^3$, $10^4$, or $3 \cdot 10^4$ optimization iterations. In addition, L-BFGS is commonly used with a relatively large learning rate, for example $\eta = 1.0$, a memory size of 100, and Wolfe line search conditions \cite{loss_landscape}. These choices reflect the fact that PINN loss landscapes can be difficult to optimize and that second-order or quasi-Newton methods may improve convergence after an Adam pre-training phase.

More recent practical recommendations suggest that the training of PINNs is highly problem-dependent and should be tuned jointly with the sampling strategy \cite{experts_guide}. For time-dependent and nonlinear PDEs such as the Allen--Cahn equation, Navier--Stokes equations, and advection-type problems, architectures with four or five hidden layers and 256 neurons per layer have been used successfully. These models are often trained for a large number of steps, for example between $10^5$ and $3 \cdot 10^5$ iterations, with learning-rate decay applied every few thousand steps. Batch sizes such as 4096 or 8192 are used for the residual loss, and the residual points are resampled at every training step rather than fixed in advance. This dynamic sampling strategy is especially relevant in higher dimensions, where a fixed set of collocation points may cover the domain poorly.

For the diffusion equation in higher dimensions, this thesis follows the general principle that the network should be large enough to approximate the solution over the full spatial domain, but not so large that optimization becomes unstable or unnecessarily expensive. A moderate baseline architecture is therefore a fully connected feed-forward network with four hidden layers and 128 neurons per layer. This is comparable to architectures used in stochastic-gradient PINN training, where Adam optimization is applied for approximately $10^5$ epochs with a linearly decaying learning rate. Larger models, such as four hidden layers with 512 neurons per layer, may be considered when the dimension is increased or when the solution contains sharper features. However, for smoother diffusion problems, smaller networks can often provide a reasonable balance between accuracy and computational cost.

The optimizer is another important hyperparameter. Adam is attractive because it is robust, easy to implement, and can handle mini-batches of residual points. A common strategy is to train with Adam for a large number of steps while decaying the learning rate, for example by multiplying it by 0.9 every $10^4$ steps. In some PINN variants, each epoch uses newly sampled residual points, such as $10^4$ points sampled from the underlying stochastic process or spatial domain. This reduces overfitting to a fixed collocation set and is particularly useful in high-dimensional domains, where full-grid discretizations are infeasible.

In comparison, L-BFGS is usually applied in full-batch mode and can be effective once the network is already close to a good solution. However, its memory cost and dependence on full residual evaluations make it less convenient for very large batches or high-dimensional sampling. For this reason, the experiments in this thesis emphasize Adam-based stochastic training, with the option of applying L-BFGS as a final refinement step when the residual batch size and memory requirements allow it.

Overall, the training setup used in this thesis is guided by the following considerations. First, smaller networks and fixed collocation sets are appropriate for classical one- and two-dimensional PINN benchmarks, as in the original PINN literature \cite{pinns_original}. Second, higher-dimensional diffusion problems benefit from wider networks, larger residual batches, and repeated resampling of collocation points. Third, Adam with learning-rate decay provides a practical baseline for long training runs, while L-BFGS can be useful as a secondary optimizer for smaller or full-batch experiments. These choices provide a compromise between approximation power, optimization stability, and computational feasibility.